\scanpage{15}
Détermination du groupe des caractères rationnels sur $\mathbb Q$

$T_{\mathbb C}=D_{\mathbb C}(M)$, $M$ réseau où $\pi_1$ opère \inserted{via un \unsure{gpe|groupe} quotient fini $\pi$ opérant fidèlement} $(\pi_1=\pi_1(\mathbb Q,\mathbb C))$

$\mathbb Z(M)$ anneau des caractères \inserted{virtuels} de $T_{\mathbb C}$, $\pi$ y opère.

$K(T)=\mathbb Z(\deleted{\unsure{N|M}}M)^\pi$\qquad Anneau de type fini sur $\mathbb Z$.

[N.B. $\mathbb Z(M)$ est l'anneau affine de $S=\deleted{\unsure{\mathbb Z}}D_{\mathbb Z}(M)$, sur lequel $\pi$ opère par automorphismes de groupes. $K(T)$ est l'anneau affine de $D_{\mathbb Z}(M)/\pi$]. Les caractères \uline{effectifs} sont ceux qui prennent des valeurs $\geq 0$ sur les points de $S(\mathbb R)^0$ ?? \margin{$\lambda_1^{n_1}\lambda_2^{n_2}\unsure{-|\cdots}\lambda_r^{n_r}$} Comment définir une structure d'effectivité \uline{motivique} ? Le procédé \unsure{standart|standard} est de \deleted{se} fixer un $i_1:G_{m\mathbb C}-T_{\mathbb C}$ i.e. un $i_1^*:M\unsure{-|\to}\mathbb Z$, avec les propriétés habituelles, \deleted{et} donnant donc $\mathbb Z(M)-\mathbb Z(\mathbb Z)$ d'où $G_{m\mathbb Z}\xrightarrow{\alpha}S$ (par \unsure{comp.|conj.} avec $\pi$....)

La condition est, \deleted{que} pour un $\chi\in K(T)\simeq\mathbb Z(S)^\pi$, que $\chi\circ\alpha\simeq\deleted{\unsure{\mathbb Z|G}}$ anneau affine de $G_{m\mathbb Z}$ soit une fonction \uline{polynômiale} i.e. \deleted{restant} finie sur $G_{m\mathbb Z}(\mathbb R)=\mathbb R^*$ pour $t\to0$.

Considérons les caractères de motifs effectifs de poids donné \unsure{$s$|$p$|$\rho$}, prouvons qu'il n'y en a qu'un nb fini. Cela \inserted{\uuline{ne}} \uuline{\unsure{tient|peut|vient}} \inserted{\uuline{pas}} simplement \unsure{du|en} \unsure{calcul|critère} \unsure{arithmétique|quantième|quand on passe de} $\mathbb Q$ par $\mathbb R$, \deleted{Alors} car si $T$ est de dim $\geq3$, la relation de finitude correspondante est \uline{fausse}.

% archive sheet number: 14
